%! Polar Function Graphs — Unit 3, Lesson 3.14 — Lesson Plan
\documentclass[10pt]{article}
\usepackage{precalculus-article}
\usepackage{precalculus-boxes}
\usepackage{graphicx}
\graphicspath{{images/}}
\newcommand{\TallMath}[1]{$\displaystyle #1\rule[-1.4em]{0pt}{3.2em}$}
\newcommand{\CourseName}{Precalculus}
\newcommand{\SchoolYear}{2026--2027}
\newcommand{\MeetingLength}{55 minutes}
\newcommand{\UnitNumberName}{Unit 3: Trigonometric and Polar Functions \quad}
\newcommand{\LessonNumberName}{Lesson 3.14: Polar Function Graphs}

\begin{document}

%----------------------------------------------------------------------
% Title Block
%----------------------------------------------------------------------
\begin{center}
  {\Large\bfseries \CourseName: \SchoolYear} \\[4pt]
  {\small\bfseries \UnitNumberName \LessonNumberName \quad (2--3 instructional periods, \MeetingLength\ each)}
\end{center}

%----------------------------------------------------------------------
% Primary Objective
%----------------------------------------------------------------------
\begin{tcolorbox}[colback=sky, colframe=navy, boxrule=0.9pt, arc=2mm,
    left=2mm, right=2mm, top=1.5mm, bottom=1.5mm]
  \textbf{Primary Objective:} Students will construct and interpret graphs of polar
  functions $r = f(\theta)$, recognizing that inputs are angle measures (from the
  positive $x$-axis) and outputs are radii (distances from the origin), and will
  describe how restricting the domain selects a portion of the polar curve.
  \textit{(Big Idea: Functions)}
\end{tcolorbox}

\vspace{0.08in}

%----------------------------------------------------------------------
% Priority Ideas & Skills
%----------------------------------------------------------------------
\begin{skillbox}[Priority Ideas \& Skills]{goldbox}
\begin{tabularx}{\linewidth}{@{} p{0.46\linewidth} X @{}}
\textbf{AP Skills} & \textbf{Key Understandings (EKs)} \\
\midrule
\textbf{2.B} \textit{Construct equivalent representations} ---
  Connect a table of $(\theta, r)$ input-output pairs to the corresponding
  polar curve, and relate the rectangular graph of $r = f(\theta)$ to the
  shape of the polar graph.
&
\textbf{EK 3.14.A.1}: The graph of $r = f(\theta)$ in polar coordinates
consists of input-output pairs where inputs are angle measures and outputs
are radii. \\[4pt]

\textbf{3.A} \textit{Describe characteristics} ---
  Describe intervals of $\theta$ for which $r > 0$, $r < 0$, or $r = 0$;
  identify maximum radius and the angles at which the curve passes through
  the pole.
&
\textbf{EK 3.14.A.2}: The domain of $r = f(\theta)$, given graphically,
can be restricted to a desired portion of the curve by selecting endpoints
corresponding to the desired angle and radius. \\[4pt]

&
\textbf{EK 3.14.A.3}: Changes in $\theta$ correspond to changes in angle
from the positive $x$-axis; changes in $r$ correspond to changes in distance
from the origin. \\
\end{tabularx}
\end{skillbox}

\vspace{0.08in}

%----------------------------------------------------------------------
% Vocabulary
%----------------------------------------------------------------------
\begin{skillbox}[{Vocabulary, Concepts, \& Theorems}]{greenbox}
\begin{tabularx}{\linewidth}{@{} p{0.22\linewidth} X @{}}
\textbf{Term} & \textbf{Definition / Note} \\
\midrule
polar function &
  A function $r = f(\theta)$ where $\theta$ is the input (angle measure from
  the positive $x$-axis) and $r$ is the output (signed distance from the origin). \\[4pt]
pole &
  The origin in polar coordinates; the point $r = 0$. The curve passes through
  the pole whenever $f(\theta) = 0$. \\[4pt]
limaçon &
  A polar curve of the form $r = a + b\sin\theta$ or $r = a + b\cos\theta$.
  When $|a| = |b|$ it is a cardioid; when $|a| < |b|$ it has an inner loop. \\[4pt]
rose curve &
  A polar curve of the form $r = a\sin(n\theta)$ or $r = a\cos(n\theta)$.
  Produces $n$ petals when $n$ is odd; $2n$ petals when $n$ is even. \\[4pt]
petal &
  One lobe of a rose curve. Traced as $r$ increases from $0$, reaches a maximum,
  then returns to $0$ over a sub-interval of $\theta$. \\[4pt]
negative $r$ &
  When $r < 0$, the point is plotted in the direction \emph{opposite} to the
  angle $\theta$ --- i.e., $180°$ from the terminal side. \\
\end{tabularx}
\end{skillbox}

\vspace{0.08in}

%----------------------------------------------------------------------
% Activate Prior Knowledge
%----------------------------------------------------------------------
\begin{skillbox}[Activate Prior Knowledge \& Spiral Review]{sky}
\begin{tabularx}{\linewidth}{@{}X c@{}}
\begin{minipage}[t]{0.96\linewidth}\vspace{0pt}
\textbf{Spiral topics activated during Warm-Up (first 8--10 minutes):}
\begin{itemize}[leftmargin=1.4em, itemsep=2pt]
  \item \textbf{Evaluating $\sin\theta$ and $\cos\theta$ at key angles}: Students need
    fluency with $\theta = 0, \pi/6, \pi/4, \pi/3, \pi/2, \pi, 3\pi/2, 2\pi$ to compute
    the output values $r = f(\theta)$ that determine where the polar curve lies.
  \item \textbf{Polar coordinate system (Lessons 3.12--3.13)}: Students must recall that
    a point $(r, \theta)$ is located distance $|r|$ from the origin along the ray at angle
    $\theta$, and that negative $r$ reverses direction.
  \item \textbf{Reading function values from a rectangular graph}: Students read $f(\theta)$
    output values at given input angles --- the same skill used to trace the polar curve
    one input-output pair at a time.
\end{itemize}
\end{minipage}
&
\end{tabularx}
\end{skillbox}

\vspace{0.08in}

%----------------------------------------------------------------------
% Hook
%----------------------------------------------------------------------
\begin{skillbox}[Hook --- Entry Event]{sky}
\textbf{Scenario:} Display a pre-drawn polar graph of $r = 2\sin\theta$ on the board.
The curve is a perfect circle of radius 1 centered at $(0, 1)$.

\textbf{Discussion prompts:}
\begin{enumerate}[leftmargin=1.4em, itemsep=3pt]
  \item \textit{``This curve is a circle --- but it comes from a sine function. How
    could a sine function possibly produce a circle?''}
  \item \textit{``Look at the table of values for $r = 2\sin\theta$.
    When $\theta = \pi/2$, what is $r$? When $\theta = 0$ or $\pi$, what is $r$?
    Plot those points mentally: do they seem to lie on a circle?''}
  \item \textit{``What happens when $\theta > \pi$? The $r$-values are negative.
    Where do those points go?''}
\end{enumerate}

\textbf{Key tension to surface:} In rectangular coordinates, a circle is not a
function. In polar coordinates, the same circle \emph{is} a function --- because
the input is an angle, not an $x$-value, and the output is a radius. The surprising
shapes of polar graphs emerge directly from reading the rectangular graph of $r = f(\theta)$
as a sequence of angle-and-radius instructions.

\textbf{Transition:} ``Today we'll read polar functions from tables and rectangular
graphs, then trace what those input-output pairs draw in the polar plane.''
\end{skillbox}

\vspace{0.08in}

%----------------------------------------------------------------------
% Lesson
%----------------------------------------------------------------------
\begin{skillbox}[Lesson --- Instruction Sequence]{sky}
\begin{multicols}{2}

\textbf{Step 1 --- What is a polar function?} (10 min)

Establish the definition: $r = f(\theta)$ assigns a radius to each angle.
Use the table for $r = 2\sin\theta$ at $\theta = 0, \pi/6, \pi/3, \pi/2, 2\pi/3,
5\pi/6, \pi$. Show each $(\theta, r)$ pair as a point on polar axes.
Emphasize EK 3.14.A.3: increasing $\theta$ rotates the ray; the $r$ value
places the point along that ray. As $\theta$ goes $0 \to \pi$, the circle
is traced. For $\theta \in (\pi, 2\pi)$, $r < 0$, so the points retrace
the same circle (opposite direction).

\columnbreak

\textbf{Step 2 --- Reading from a rectangular graph} (12 min)

Display the rectangular graph of $r = \sin(2\theta)$ for $\theta \in [0, 2\pi]$.
Walk through: where is $r = 0$? (at $\theta = 0, \pi/2, \pi, 3\pi/2, 2\pi$)
Where is $r$ at its maximum of $1$? (at $\theta = \pi/4, 5\pi/4$)
Where is $r$ at its minimum of $-1$? (at $\theta = 3\pi/4, 7\pi/4$)
Each zero corresponds to the curve passing through the pole.
Each maximum or minimum corresponds to the tip of a petal.
Connect: 4 zeros \emph{between} $0$ and $2\pi$ on the rectangular graph
$\Rightarrow$ 4 petals on the polar graph.

\end{multicols}

\begin{multicols}{2}

\textbf{Step 3 --- Families of polar curves} (10 min)

Limaçons $r = a + b\sin\theta$: trace using the table from the activity
(e.g., $r = 1 + \cos\theta$). Max $r = a + b$ at $\theta = 0$; $r = 0$
when $a + b\cos\theta = 0$. Circles: $r = a$ is a circle of radius $|a|$
centered at the origin (constant output).
Rose curves $r = a\sin(n\theta)$: odd $n \Rightarrow n$ petals;
even $n \Rightarrow 2n$ petals.

\columnbreak

\textbf{Step 4 --- Domain restriction} (8 min)

Demonstrate EK 3.14.A.2 using the polar graph of $r = \sin(2\theta)$.
Ask: if we only allow $\theta \in [0, \pi/2]$, which portion of the curve
is traced? (Just the first petal, in the first quadrant.)
If $\theta \in [0, \pi]$, two petals. The domain restriction selects a
sub-arc of the full curve.

\end{multicols}

\smallskip
\textbf{Pacing note:} Steps 1--2 in Period 1; Steps 3--4 and group activity
in Period 2; review and exit ticket in Period 3 if needed.
\end{skillbox}

\vspace{0.08in}

%----------------------------------------------------------------------
% Active Monitoring
%----------------------------------------------------------------------
\begin{skillbox}[Active Monitoring]{redbox}
\begin{itemize}[leftmargin=1.4em, itemsep=3pt]
  \item \textbf{Confusing $\theta$ and $r$ roles}: Students write $r$ on the horizontal
    axis and $\theta$ on the vertical. Prompt: \textit{``In polar coordinates, which
    quantity is the angle and which is the distance from the origin?''}
  \item \textbf{Ignoring negative $r$ values}: Students skip $\theta > \pi$ for
    $r = 2\sin\theta$ because $r < 0$. Prompt: \textit{``Negative $r$ means the
    point is in the opposite direction. Where does $(\theta = 3\pi/2,\ r = -2)$ plot?''}
  \item \textbf{Petal count error for even $n$}: Students predict $n$ petals for
    $r = \sin(2\theta)$. Counter: \textit{``Count the zeros of the rectangular graph
    in $[0, 2\pi]$. Each zero-to-zero interval traces one petal --- how many intervals
    are there?''}
  \item \textbf{Domain restriction confusion}: Students think restricting the domain
    removes petals permanently. Clarify: \textit{``The full curve is still defined;
    we are only tracing part of it.''}
\end{itemize}
\end{skillbox}

\vspace{0.08in}

%----------------------------------------------------------------------
% Group Work & Differentiation
%----------------------------------------------------------------------
\begin{skillbox}[Group Work \& Differentiation]{redbox}
Students work in groups of 2--3 on the tiered activity (teacher-assigned
or student self-selected with guidance).

\begin{itemize}[leftmargin=1.4em, itemsep=4pt]
  \item \textbf{Tier R (Reinforce)}: Given a complete input-output table for
    $r = 1 + \cos\theta$ (12 angle values), students answer comprehension
    questions reading the table: when is $r = 0$? When is $r$ maximum?
    What is the maximum value of $r$?
  \item \textbf{Tier A (Apply)}: Given a pre-drawn rectangular graph of
    $r = \sin(2\theta)$ for $\theta \in [0, 2\pi]$, students identify
    intervals where $r$ is positive, negative, or zero, and determine the
    petal count of the polar curve with written justification.
  \item \textbf{Tier E (Extend)}: Given pre-drawn polar graphs of
    $r = 2\sin\theta$ and $r = 1 + \sin\theta$, students compare maximum
    $r$, $\theta$ values where $r = 0$, and describe in words how each
    curve is traced as $\theta$ goes from $0$ to $2\pi$.
\end{itemize}
\end{skillbox}

\vspace{0.08in}

%----------------------------------------------------------------------
% Individual Work & Assessment
%----------------------------------------------------------------------
\begin{skillbox}[Individual Work \& Assessment]{redbox}
Exit ticket administered independently (final 5--7 minutes of Period 2 or 3).

\begin{enumerate}[leftmargin=1.4em, itemsep=4pt]
  \item Students complete a short table of $r = 3\cos\theta$ at
    $\theta = 0, \pi/2, \pi, 3\pi/2$.
  \item Students identify the $\theta$ values (in $[0, 2\pi)$) where $r = 0$.
  \item Multiple-choice: identify which statement correctly describes the
    input/output roles of $\theta$ and $r$ in $r = f(\theta)$.
\end{enumerate}

\textbf{Data use:} Sort exit tickets for role-confusion errors (Q3) and
computation errors (Q1). Students who confuse inputs and outputs need
a direct re-teach of the polar function definition at the start of
Lesson 3.15.
\end{skillbox}

\vspace{0.08in}

%----------------------------------------------------------------------
% Reinforcement & Extension
%----------------------------------------------------------------------
\begin{skillbox}[Reinforcement \& Extension]{goldbox}
\textbf{Homework (Lesson 3.14):} 5--6 problems: read a table for $r = 1 - \sin\theta$
and answer questions about maximum $r$ and zeros; interpret a rectangular graph of
$r = \cos(3\theta)$ to predict the petal count of the polar curve; domain restriction
problem; 2 AP-style MC items. No sketching.

\textbf{Preview of Lesson 3.15 --- Rates of Change of Polar Functions:}
Ask students: \textit{``Now that you can read a polar graph from a table or rectangular
graph, what does it mean for the radius to be \emph{changing} at a given angle?
In Lesson 3.15 we'll quantify how fast $r$ is changing with respect to $\theta$,
and connect that to the direction the curve is moving.''}
\end{skillbox}

\end{document}
